% \section{Problem Setup: Sparse Sensors, Wind, and Source Apportionment}
% \label{sec:problem_setup}

% \paragraph{City-Scale Pollution as a Transport System}

% We consider city-scale air pollution as a spatio-temporal transport
% process over a spatial domain \(\Omega \subset \mathbb{R}^d\). Let
% \(u(\mathbf x,t)\in \mathbb R\) denote the pollutant concentration
% field at location \(\mathbf x\) and time \(t\). The evolution of
% \(u\) is governed by a combination of transport, dispersion, source
% emissions, and background variation. A generic continuous-time
% description is
% \begin{equation}
% \frac{\partial u}{\partial t}
% =
% \mathcal F_{\eta}
% \left(
% u,\mathbf x,t;\mathbf w(t)
% \right)
% +
% s(\mathbf x,t)
% +
% b(\mathbf x,t),
% \label{eq:generic_pollution_dynamics}
% \end{equation}
% where \(\mathbf w(t)\) denotes the meteorological state, including wind
% speed and direction, \(s(\mathbf x,t)\) denotes source emissions,
% \(b(\mathbf x,t)\) denotes background or regional pollution, and
% \(\eta\) represents physical parameters such as dispersion rates,
% mixing assumptions, or plume-model parameters.

% In practice, the exact dynamics \(\mathcal F_{\eta}\) are only
% partially known. City-scale transport models rely on simplified physics,
% coarse meteorological inputs, approximate boundary conditions, and
% imperfect source inventories. Thus, source apportionment from sparse
% sensors is not simply a supervised prediction problem: it is an inverse
% problem with sparse concentration observations, sparse wind observations,
% partial physical knowledge, finite transport delays, and background
% temporal variation.

% \paragraph{Spatial Discretization and Sparse Observation Operator}

% We discretize the concentration field on \(n\) spatial grid cells, so
% that the concentration field at time \(t\) is represented by
% \begin{equation}
% \mathbf u_t \in \mathbb R^n .
% \end{equation}
% The full concentration field is not observed. Instead, measurements are
% available at a small number of fixed monitoring stations. Let
% \[
% X_{\mathrm{sens}}
% =
% \{\mathbf x_1,\ldots,\mathbf x_m\}
% \]
% denote the sensor locations, with \(m \ll n\). The observation operator
% \(\mathcal O\) maps the latent concentration field to sensor readings:
% \begin{equation}
% \mathbf y_t
% =
% \mathcal O(\mathbf u_t)
% +
% \boldsymbol\epsilon_t
% \in \mathbb R^m,
% \label{eq:sparse_sensor_observation}
% \end{equation}
% where \(\boldsymbol\epsilon_t\) denotes sensor noise. When sensors are
% located on grid cells, \(\mathcal O\) selects the corresponding entries
% of \(\mathbf u_t\); more generally, \(\mathcal O\) may interpolate the
% field to sensor locations.

% Stacking observations over a finite horizon \(t=1,\ldots,T\), we obtain
% \begin{equation}
% \mathbf Y
% =
% \begin{bmatrix}
% \mathbf y_1\\
% \vdots\\
% \mathbf y_T
% \end{bmatrix}
% \in
% \mathbb R^{mT}.
% \label{eq:stacked_sensor_observations}
% \end{equation}
% Since \(m\) is typically much smaller than the number of spatial grid
% cells, the observation operator is highly non-invertible. Consequently,
% many distinct latent concentration fields, source configurations, and
% background patterns may be consistent with the same sparse sensor
% trajectory.

% \paragraph{Known Source Inventories}

% Rather than attempting to recover an arbitrary fine-scale unknown source
% field, we focus on inventory-based source apportionment. We assume that
% a set of \(K\) known or proxy source maps is available:
% \begin{equation}
% S
% =
% \begin{bmatrix}
% \mathbf s_1 & \mathbf s_2 & \cdots & \mathbf s_K
% \end{bmatrix}
% \in
% \mathbb R^{q\times K},
% \label{eq:source_inventory_matrix}
% \end{equation}
% where each column \(\mathbf s_k\in\mathbb R^q\) represents the spatial
% emission pattern of a source group, such as traffic, industrial activity,
% brick kilns, population-related emissions, road dust, or other inventory
% categories. Here \(q\) denotes the number of source-grid cells. Often
% \(q=n\), but we keep the notation separate because source inventories
% and concentration fields may be represented at different spatial
% resolutions.

% The goal of source apportionment is to estimate the contribution of each
% source group to the observed concentration signal. Let
% \(\boldsymbol\theta\in\mathbb R^K\) denote the source contribution or
% activity vector. In the simplest setting, \(\boldsymbol\theta\) is assumed
% constant over the observation window, with
% \begin{equation}
% \boldsymbol\theta \ge 0.
% \end{equation}
% More generally, source activity may vary over time. We begin with the
% constant-window case and discuss time-varying extensions later.

% \paragraph{Lagged Wind-Conditioned Transport Response}

% Meteorology determines not only the direction and spread of pollutant
% transport, but also the delay with which emissions become visible at the
% sensor network. Emissions released at time \(t-\ell\) may contribute to
% sensor readings at time \(t\) only after being transported by the
% intervening wind field. Thus, the source-to-sensor response is generally
% not instantaneous.

% Let \(L\) denote the maximum transport lag considered over the observation
% window. For each sensor time \(t\) and lag \(\ell\), we define a lagged
% transport operator
% \begin{equation}
% G_{t,t-\ell}
% \left(
% \mathbf w_{t-\ell:t}
% \right)
% \in
% \mathbb R^{n\times q},
% \label{eq:lagged_transport_operator}
% \end{equation}
% which maps emissions released on the source grid at time \(t-\ell\) to
% their induced concentration field on the concentration grid at time \(t\),
% after transport under the wind sequence \(\mathbf w_{t-\ell:t}\). This
% operator may be constructed from a time-stepping advection--diffusion
% solver, a Gaussian plume model with travel-time correction, a
% wind-shifted kernel, or a learned surrogate constrained by physical
% structure.

% For times near the beginning of the observation window, only lags that
% fall within the available history are included. Define
% \begin{equation}
% \mathcal L_t
% =
% \{0,1,\ldots,\min(L,t-1)\}.
% \label{eq:available_lag_set}
% \end{equation}
% The lagged source-response matrix at sensor time \(t\) is then
% \begin{equation}
% H_t^{\mathrm{lag}}
% :=
% \sum_{\ell\in \mathcal L_t}
% \mathcal O
% G_{t,t-\ell}
% \left(
% \mathbf w_{t-\ell:t}
% \right)
% S
% \in
% \mathbb R^{m\times K}.
% \label{eq:lagged_source_response_matrix}
% \end{equation}
% The \(k\)-th column of \(H_t^{\mathrm{lag}}\) is the cumulative
% sensor-space fingerprint of source group \(k\) at time \(t\), accounting
% for emissions from recent past times that can reach the sensors within
% the transport-lag window.

% This lagged formulation includes the instantaneous model as a special
% case. When \(L=0\), we recover
% \begin{equation}
% H_t^{\mathrm{lag}}
% =
% \mathcal O G_{t,t}(\mathbf w_t)S.
% \end{equation}

% \paragraph{Finite-Horizon Lagged Source Apportionment Model}

% After background and temporal correction, the sensor observation at time
% \(t\) is modeled as
% \begin{equation}
% \tilde{\mathbf y}_t
% =
% H_t^{\mathrm{lag}}\boldsymbol\theta
% +
% \boldsymbol\epsilon_t,
% \label{eq:single_time_lagged_apportionment}
% \end{equation}
% where \(\tilde{\mathbf y}_t\) denotes the corrected concentration signal.
% Stacking over time gives
% \begin{equation}
% \tilde{\mathbf Y}
% =
% H_{1:T}^{\mathrm{lag}}\boldsymbol\theta
% +
% \mathbf E,
% \label{eq:stacked_lagged_apportionment_model}
% \end{equation}
% where
% \begin{equation}
% H_{1:T}^{\mathrm{lag}}
% =
% \begin{bmatrix}
% H_1^{\mathrm{lag}}\\
% H_2^{\mathrm{lag}}\\
% \vdots\\
% H_T^{\mathrm{lag}}
% \end{bmatrix}
% \in
% \mathbb R^{mT\times K}.
% \label{eq:stacked_lagged_source_response_matrix}
% \end{equation}

% The matrix \(H_{1:T}^{\mathrm{lag}}\) is the central object in our
% analysis. It encodes how known source groups appear at the sparse sensor
% network under the observed wind sequence, after accounting for finite
% transport delays. Source groups are distinguishable only if their
% lagged wind-conditioned fingerprints are sufficiently independent in
% stacked sensor space.

% \paragraph{Background and Temporal Confounding}

% Real air-quality observations include background pollution, regional
% transport, time-of-day effects, seasonal variation, low-frequency drift,
% and sensor offsets that are not explained by the local source inventory.
% We model these effects using a low-dimensional background-correction
% basis
% \begin{equation}
% Q\in\mathbb R^{mT\times r},
% \end{equation}
% with coefficients \(\boldsymbol\gamma\in\mathbb R^r\). The observation
% model becomes
% \begin{equation}
% \mathbf Y
% =
% H_{1:T}^{\mathrm{lag}}\boldsymbol\theta
% +
% Q\boldsymbol\gamma
% +
% \mathbf E.
% \label{eq:lagged_apportionment_with_background}
% \end{equation}

% This background component is important for identifiability. A source
% group may appear distinguishable before background correction, but become
% confounded once time-of-day or regional background variation is allowed
% to explain part of the signal. To account for this, we project the data
% and source-response matrix onto the orthogonal complement of the
% background-correction space. Let
% \begin{equation}
% P_Q^\perp
% =
% I
% -
% Q(Q^\top Q)^{-1}Q^\top
% \label{eq:background_projection}
% \end{equation}
% denote the background-removal projection. If \(Q^\top Q\) is singular or
% ill-conditioned, the inverse can be replaced with a pseudoinverse. The
% corrected inverse problem is
% \begin{equation}
% P_Q^\perp \mathbf Y
% =
% P_Q^\perp H_{1:T}^{\mathrm{lag}}\boldsymbol\theta
% +
% P_Q^\perp\mathbf E.
% \label{eq:projected_lagged_apportionment_model}
% \end{equation}
% Thus, source apportionment should be analyzed using the projected lagged
% response matrix \(P_Q^\perp H_{1:T}^{\mathrm{lag}}\), rather than the raw
% response matrix \(H_{1:T}^{\mathrm{lag}}\).

% \paragraph{Lagged Source Fingerprints and Ambiguity}

% Let
% \begin{equation}
% \mathbf h_k^{\mathrm{lag}}
% =
% \begin{bmatrix}
% \displaystyle
% \sum_{\ell\in \mathcal L_1}
% \mathcal O
% G_{1,1-\ell}
% \left(
% \mathbf w_{1-\ell:1}
% \right)
% \mathbf s_k\\[1.2em]
% \displaystyle
% \sum_{\ell\in \mathcal L_2}
% \mathcal O
% G_{2,2-\ell}
% \left(
% \mathbf w_{2-\ell:2}
% \right)
% \mathbf s_k\\
% \vdots\\
% \displaystyle
% \sum_{\ell\in \mathcal L_T}
% \mathcal O
% G_{T,T-\ell}
% \left(
% \mathbf w_{T-\ell:T}
% \right)
% \mathbf s_k
% \end{bmatrix}
% \in
% \mathbb R^{mT}
% \label{eq:lagged_source_fingerprint}
% \end{equation}
% denote the finite-horizon lagged sensor fingerprint of source group \(k\).
% The columns of \(H_{1:T}^{\mathrm{lag}}\) are precisely these
% fingerprints:
% \begin{equation}
% H_{1:T}^{\mathrm{lag}}
% =
% \begin{bmatrix}
% \mathbf h_1^{\mathrm{lag}} &
% \mathbf h_2^{\mathrm{lag}} &
% \cdots &
% \mathbf h_K^{\mathrm{lag}}
% \end{bmatrix}.
% \end{equation}

% Two kinds of ambiguity arise naturally. First, a source group may be
% weakly visible if its projected fingerprint has small norm:
% \begin{equation}
% \left\|
% P_Q^\perp
% \mathbf h_k^{\mathrm{lag}}
% \right\|
% \approx 0.
% \end{equation}
% Such a source has little measurable effect on the sensor network under
% the available wind regime and background correction. Second, two source
% groups may be difficult to separate even if both are visible, if their
% projected fingerprints are nearly proportional:
% \begin{equation}
% P_Q^\perp \mathbf h_i^{\mathrm{lag}}
% \approx
% c\,
% P_Q^\perp \mathbf h_j^{\mathrm{lag}}
% \quad
% \text{for some } c>0.
% \end{equation}
% This captures the amplitude--distance ambiguity in source apportionment:
% a strong distant source and a weak nearby source may produce similar
% sensor readings after transport attenuation and background correction.

% We quantify pairwise source ambiguity using projected fingerprint
% coherence:
% \begin{equation}
% \rho_{ij}
% =
% \frac{
% \left|
% \left(
% P_Q^\perp \mathbf h_i^{\mathrm{lag}}
% \right)^\top
% \left(
% P_Q^\perp \mathbf h_j^{\mathrm{lag}}
% \right)
% \right|
% }{
% \left\|
% P_Q^\perp \mathbf h_i^{\mathrm{lag}}
% \right\|
% \left\|
% P_Q^\perp \mathbf h_j^{\mathrm{lag}}
% \right\|
% }.
% \label{eq:projected_lagged_source_fingerprint_coherence}
% \end{equation}
% Values \(\rho_{ij}\approx 1\) indicate that source groups \(i\) and \(j\)
% are nearly indistinguishable up to a scaling factor. In practice, such
% source groups should not be interpreted separately without additional
% sensing, stronger priors, or a coarser source grouping.

% \paragraph{Identifiability Objective}

% The goal of this work is not merely to estimate source contributions,
% but to determine when such estimates are identifiable from the available
% sensors, wind observations, source inventories, background correction,
% and transport-lag model. In the projected lagged model
% \eqref{eq:projected_lagged_apportionment_model}, exact identifiability of
% \(\boldsymbol\theta\) requires
% \begin{equation}
% \operatorname{rank}
% \left(
% P_Q^\perp H_{1:T}^{\mathrm{lag}}
% \right)
% =
% K.
% \label{eq:lagged_rank_identifiability_condition}
% \end{equation}
% Noise-robust identifiability further depends on the conditioning of
% \(P_Q^\perp H_{1:T}^{\mathrm{lag}}\), especially its smallest singular
% value:
% \begin{equation}
% \sigma_{\min}
% \left(
% P_Q^\perp H_{1:T}^{\mathrm{lag}}
% \right).
% \label{eq:lagged_sigma_min_identifiability}
% \end{equation}

% This perspective reframes sparse-sensor source apportionment as an
% identifiability problem. Known inventories make the problem meaningful,
% but they do not guarantee separability. Source contributions can be
% reliably estimated only when lagged wind-conditioned source fingerprints
% remain sufficiently independent after accounting for sparse sensing, wind
% uncertainty, background variation, model approximation, and finite
% transport delays.


\section{Problem Setup: Sparse Sensors, Wind, and Source Apportionment}
\label{sec:problem_setup}

We formalize source apportionment as a finite-dimensional inverse problem in
which known source inventories are transported to a sparse sensor network under
a wind-conditioned lagged operator.  This section defines the observation
model, the inventory parameterization, the background correction, and the
projected response matrix used throughout the paper.

\paragraph{City-scale transport model.}
Let \(u(\mathbf x,t)\in \mathbb R\) denote the pollutant concentration field on
\(\Omega\subset \mathbb R^d\).  A generic continuous-time model is
\begin{equation}
\frac{\partial u}{\partial t}
= \mathcal F_\eta(u,\mathbf x,t;w(t)) + s(\mathbf x,t)+b(\mathbf x,t),
\label{eq:continuous_transport}
\end{equation}
where \(w(t)\) is the meteorological state, \(s\) is local source emission,
\(b\) is background or regional pollution, and \(\eta\) denotes transport and
dispersion parameters.  The exact dynamics may be approximated by a plume
model, an advection--diffusion solver, or a learned surrogate constrained by
transport structure.

\paragraph{Sparse observations.}
Discretize the concentration field on \(n\) grid cells, so that
\(\mathbf u_t\in \mathbb R^n\).  Let
\(X_{\mathrm{sens}}=\{\mathbf x_1,\ldots,\mathbf x_m\}\) be the fixed sensor
locations with \(m\ll n\).  The observation operator
\(O:\mathbb R^n\to \mathbb R^m\) samples or interpolates the latent field at
these locations:
\begin{equation}
\mathbf y_t = O\mathbf u_t + \boldsymbol\epsilon_t,
\qquad \mathbf y_t\in \mathbb R^m .
\label{eq:sensor_obs}
\end{equation}
Stacking over \(t=1,\ldots,T\) gives
\begin{equation}
Y = [\mathbf y_1^\top,\ldots,\mathbf y_T^\top]^\top \in \mathbb R^{mT}.
\label{eq:stacked_observations}
\end{equation}
Because \(m\) is small relative to the field dimension, many latent fields and
source configurations can be consistent with the same sensor trajectory.

\paragraph{Inventory-based source parameterization.}
We assume a set of \(K\) known or
proxy source maps,
\begin{equation}
S = [\mathbf s_1\;\mathbf s_2\;\cdots\;\mathbf s_K]\in \mathbb R^{q\times K},
\label{eq:inventory_matrix}
\end{equation}
where \(\mathbf s_k\in\mathbb R^q\) is the spatial pattern of source group
\(k\).  The source grid dimension \(q\) may differ from the concentration-grid
dimension \(n\).  We represent time-varying source activity with a shared
low-dimensional temporal basis
\begin{equation}
\Phi=[\boldsymbol\phi_1\;\cdots\;\boldsymbol\phi_B]
\in\mathbb R_+^{T\times B},
\qquad
C=[c_{kb}]\in\mathbb R_+^{K\times B}.
\label{eq:temporal_activity_basis}
\end{equation}
The activity of source \(k\) at time \(t\) is
\begin{equation}
\theta_k(t)=\sum_{b=1}^{B}c_{kb}\phi_b(t),
\qquad \theta_k(t)\ge0.
\label{eq:source_activity}
\end{equation}
The activity bases used for source emissions are nonnegative; signed temporal
harmonics, when used, belong to the separate background basis \(Q\).  We
vectorize \(C\) in source-major, basis-minor order as
\(\mathbf c=\operatorname{vec}_{\mathrm{src}}(C)\in\mathbb R_+^J\), where
\(J=KB\).  The constant-window model is the special case \(B=1\),
\(\phi_1(t)=1\), and \(c_{k1}=\theta_k\).

\paragraph{Lagged wind-conditioned response.}
Emissions released at time \(t-\ell\) may affect sensors at time \(t\) only
after being transported by the intervening wind field.  Let
\begin{equation}
G_{t,t-\ell}^{\partial}(w_{t-\ell:t})\in \mathbb R^{n\times q}
\label{eq:lagged_operator}
\end{equation}
map emissions released on the source grid at time \(t-\ell\) to the
concentration grid at time \(t\).  The superscript \(\partial\) denotes an
open-boundary operator: mass transported outside \(\Omega\) leaves the modeled
domain rather than being reflected or wrapped around.  For a maximum lag \(L\),
define
\begin{equation}
\mathcal L_t = \{0,1,\ldots,\min(L,t-1)\}.
\label{eq:lag_set}
\end{equation}
For source \(k\) and basis component \(b\), the corresponding sensor response
at time \(t\) is
\begin{equation}
\mathbf h_{t,kb}^{\mathrm{lag}}
= \sum_{\ell\in\mathcal L_t}
\phi_b(t-\ell)
O G_{t,t-\ell}^{\partial}(w_{t-\ell:t})\mathbf s_k
\in \mathbb R^m.
\label{eq:time_response}
\end{equation}
Stacking over time and arranging columns in source-major, basis-minor order
gives
\begin{equation}
\begin{aligned}
H_{\Phi}^{\mathrm{lag}}
&=
[\mathbf h_{11}^{\mathrm{lag}}\;\cdots\;
 \mathbf h_{1B}^{\mathrm{lag}}\;\cdots\;
 \mathbf h_{KB}^{\mathrm{lag}}]
\in \mathbb R^{mT\times J},\\
J&=KB,
\end{aligned}
\label{eq:stacked_response}
\end{equation}
where each \(\mathbf h_{kb}^{\mathrm{lag}}\in\mathbb R^{mT}\) is the
finite-horizon source--basis, or coefficient, fingerprint.  When \(B=1\),
these columns reduce to the source fingerprints of the constant-activity model.

\paragraph{Background and temporal correction.}
Observed concentrations include components not explained by the local source
inventory, including regional background, smooth temporal trends, and
sensor-specific offsets.  We represent these components with a low-dimensional
basis
\begin{equation}
Q\in \mathbb R^{mT\times r}, \qquad \boldsymbol\gamma\in \mathbb R^r,
\label{eq:background_basis}
\end{equation}
and write
\begin{equation}
Y = H_{\Phi}^{\mathrm{lag}}\mathbf c + Q\boldsymbol\gamma + E.
\label{eq:model_with_background}
\end{equation}
Let
\begin{equation}
P_Q^\perp = I - QQ^\dagger
\label{eq:background_projection}
\end{equation}
be the projection onto the orthogonal complement of the background space, with
\(Q^\dagger\) denoting the Moore--Penrose pseudoinverse.  The corrected inverse
problem is
\begin{equation}
\widetilde Y = \widetilde H_\Phi\mathbf c + \widetilde E,
\qquad
\widetilde Y=P_Q^\perp Y,
\qquad
\widetilde H_\Phi=P_Q^\perp H_{\Phi}^{\mathrm{lag}}.
\label{eq:projected_model}
\end{equation}
All identifiability diagnostics in this paper are computed from
\(\widetilde H_\Phi\), not from the raw inventory matrix or the unprojected
response.

\paragraph{Projected source--basis fingerprints.}
Writing
\begin{equation}
H_{\Phi}^{\mathrm{lag}}
=[\mathbf h_{11}^{\mathrm{lag}}\;\cdots\;\mathbf h_{KB}^{\mathrm{lag}}],
\qquad
\widetilde H_\Phi
=[\widetilde{\mathbf h}_{11}\;\cdots\;\widetilde{\mathbf h}_{KB}],
\label{eq:source_fingerprints}
\end{equation}
we have
\(\widetilde{\mathbf h}_{kb}=P_Q^\perp\mathbf h_{kb}^{\mathrm{lag}}\).
These projected fingerprints encode the observable effect of each temporal
component of each source after wind transport, lag, sparse sensing, and
background correction.  Source apportionment at the chosen temporal-basis
resolution is well posed only to the extent that these coefficient fingerprints
are visible and separable.
